\documentclass[conference,transmag,onecolumn,11pt]{IEEEtran}
\usepackage{include/atcatg}
\usepackage{wrapfig}
\setcounter{figure}{0}

\begin{document}

% PAPER TITLE
% Titles are generally capitalized except for words such as a, an, and, as,
% at, but, by, for, in, nor, of, on, or, the, to and up, which are usually
% not capitalized unless they are the first or last word of the title.
% Linebreaks \\ can be used within to get better formatting as desired.
% Do not put math or special symbols in the title except for chemical compositions.

\title{Presentation Title Here, Chemical Example ABO\textsubscript{3}}

% AUTHOR NAMES AND AFFILIATIONS
% Author names should be given in Latin script (English) with the given name formatted normally,
% followed by family/patronymic/surname in the \familyname{YOURNAME} block.
% If you cannot input accented letters from your keyboard, please see the following website for accent codes: https://www.overleaf.com/learn/latex/International_language_support#Reference_guide
% The presenter should be noted in the \presenter[email_address]{FULL NAME HERE} block, preferably with an email address where conference attendees can contact you included
% Follow each name with \affiliations{number} - authors can have multiple affiliations separated by commas
% Each line of authors should be on a separate \authorline{} - do not use \\ to force line breaks
% Follow all names with affiliations, \affiliation{number}{Institution Name, Country}
% Please do not distinguish between different departments at the same institution, or include full addresses, but do include the country name

\author{
    \authorline{  \presenter[email@domain.tld ]{First \familyname{Author}}\affiliations{1},
        Second \familyname{Author}\affiliations{1},
        Third \familyname{Author}\affiliations{2}
        Fourth \familyname{Author}\affiliations{2}
        }
    \affiliation{1}{Institution One Name, Country}
    \affiliation{2}{Institution Two Name, Country}
}

\maketitle
%\vspace{-4em}

%FIGURE AND CAPTION
%This function includes the figure in the abstract, currently a test so may need reformatting.
% [width = 0.7\linewidth] will likely need adjusting on a case-by-case basis, but can be done towards the end.
%The figure should be saved in the "graphics\abstractfigures" subfolder and labelled similarly to the abstract file, i.e. REGION_FirstName_LastName.imagetype
%The caption should be included in the \caption{} environment as plain text.


\begin{wrapfigure}{R}{0.4\textwidth}
\includegraphics[width=0.4\textwidth]{graphics/abstractfigures/Region_FirstName_LastName.png}
\caption{IEEE Atc Atg conference}
\label{figure1}
\end{wrapfigure}

% ABSTRACT TEXT
% Please put the full abstract in English here.
% Text and inline equations only, no figures, tables or block equations. 
% The full length of the abstract text must fit within the text box.

\abstracttext{Please type your abstract in English here. \textbf{The maximum number of words is 150}. Longer abstracts will not be considered. The abstract can include text, mathematical symbols, and inline equations (for example $T_{C}=500K$, or $\mu_{B}=-\frac{e\hbar}{2mec}$) only. 
You can refer to pre-prints \cite{ref1}, ideally providing a link for online access \cite{ref2, ref3}. The \textbf{total number of references should not exceed 3}. Please refer to the \texttt{README.txt} for submission instructions.}

\vspace{3em}

\section*{Acknowledgements}
Please type here the acknowledgements to your work. \textbf{The maximum number of words is 150}.


\vspace{3em}

%Put the bibliography in a separate file named Region_FirstName_LastName.bib and include it in the same folder as this document

\bibliographystyle{include/IEEEtran-notitle}
\bibliography{Region_FirstName_LastName}

\end{document}